% =====================================================================
%  1bat-ud5-10.tex  —  SESSIÓ 10: Actualitzem el lloguer: IPC i índexs
%  (sense preàmbul: s'inclou des de main.tex)
% =====================================================================
\horatitol{Sessió 10 — Actualitzem el lloguer: IPC i índexs}%
{El propietari apuja el lloguer segons l'IPC: calculem la nova quota i aprenem a llegir una taula de nombres índex i taxes de variació.}

\subsection{Idea clau}

Tornem als percentatges de les primeres sessions, ara en un context plenament real i
proper a la teva feina futura. Una família usuària dels serveis socials paga un
lloguer de $\num{650}\eur$ i el propietari l'apuja segons l'\textbf{IPC interanual}.
Avui calculem la nova quota i, sobretot, aprenem a llegir una \textbf{taula d'IPC}:
què és un \textbf{nombre índex} i una \textbf{taxa de variació}.

\begin{idea}
\textbf{IPC, nombres índex i taxa de variació}
\begin{itemize}[leftmargin=1.4em, itemsep=2pt]
  \item \textbf{IPC} (índex de preus de consum): mesura com evolucionen, de mitjana,
        els preus dels productes que consumeix una família.
  \item \textbf{Nombre índex:} un valor referit a una base $=100$. Si l'IPC passa de
        $100$ a $108$, els preus han pujat un $8\%$ respecte de la base.
  \item \textbf{Taxa de variació interanual:} el percentatge que han pujat (o baixat)
        els preus d'un any respecte de l'anterior.
  \item \textbf{Actualitzar amb l'IPC:} apujar una quantitat amb l'índex de variació
        corresponent. Pujar un $3,5\%$ és multiplicar per $1,035$.
\end{itemize}
\end{idea}

\subsection{Exemples}

\begin{exemple}[la nova quota de lloguer]
El lloguer és de $\num{650}\eur$ i l'IPC interanual és del $3,5\%$. La nova quota
s'obté aplicant l'índex de variació $1,035$:
\[
\num{650}\per 1,035=\num{672,75}\eur.
\]
L'augment mensual és $\num{672,75}-\num{650}=\num{22,75}\eur$. Sembla poc, però al cap
d'un any són $12\per\num{22,75}=\num{273}\eur$ més que la família ha de pagar només de
lloguer. Una pujada de l'IPC no és una xifra abstracta: són \textbf{diners reals} cada
mes.
\end{exemple}

\begin{exemple}[llegir una taula d'IPC]
Aquesta taula mostra l'evolució d'un IPC amb base $100$ l'any $2021$:
\begin{center}
\renewcommand{\arraystretch}{1.2}
\begin{tabular}{c c c}
\toprule
\textbf{Any} & \textbf{Índex (base 2021)} & \textbf{Variació interanual} \\
\midrule
$2021$ & $100,0$ & — \\
$2022$ & $108,4$ & $+8,4\%$ \\
$2023$ & $114,2$ & $+5,4\%$ \\
$2024$ & $117,9$ & $+3,2\%$ \\
\bottomrule
\end{tabular}
\end{center}
La columna d'índex es llegeix respecte de la base: el $2024$, l'índex $117,9$ vol dir
que els preus han pujat un $17,9\%$ \emph{des del $2021$}. La variació
\textbf{interanual} compara cada any amb l'anterior: del $2023$ al $2024$, per exemple,
$\dfrac{117,9}{114,2}\approx 1,032$, és a dir un $+3,2\%$.
\end{exemple}

\subsection{Exercicis resolts}

\begin{exresolt}[actualitzar una quota amb l'IPC]
Un lloguer de $\num{800}\eur$ s'actualitza amb un IPC del $2,8\%$. Calcula la nova
quota i l'augment anual que suposa.
\solucio
Índex de variació: $1+0,028=1,028$. Nova quota:
\[
\num{800}\per 1,028=\num{822,4}\eur.
\]
Augment mensual: $\num{822,4}-\num{800}=\num{22,4}\eur$; augment anual:
$12\per\num{22,4}=\num{268,8}\eur$.
\resposta{Nova quota $\num{822,4}\eur$; $\num{268,8}\eur$ més a l'any.}
\end{exresolt}

\begin{exresolt}[interpretar nombres índex]
Si l'IPC d'alimentació passa de $100$ (base) a $112$ en dos anys, quant han pujat els
preus dels aliments? I si un producte costava $\num{4}\eur$ a la base, quant costa
ara, de mitjana?
\solucio
L'índex $112$ sobre base $100$ vol dir que els preus han pujat un
$\dfrac{112-100}{100}=12\%$.

Un producte de $\num{4}\eur$ a la base costa ara, de mitjana,
$\num{4}\per 1,12=\num{4,48}\eur$.
\resposta{Pujada del $12\%$; el producte de $\num{4}\eur$ passa a $\num{4,48}\eur$.}
\end{exresolt}

\subsection{Exercicis per fer}

\begin{exercicis}
\begin{exlist}
  \item Un lloguer de $\num{700}\eur$ s'actualitza amb un IPC del $4\%$. Calcula la
        nova quota i l'augment anual.
  \item Un IPC passa de $100$ (base) a $106,5$. Quant han pujat els preus respecte de
        la base? Un producte de $\num{20}\eur$ a la base, quant costa ara?
  \item De la taula d'IPC de l'exemple, calcula la variació interanual del $2021$ al
        $2022$ a partir dels índexs ($100$ i $108,4$) i comprova que dóna el $+8,4\%$
        indicat.
  \item Una pensió de $\num{1100}\eur$ s'actualitza amb un IPC del $3\%$. Quant cobrarà
        la persona després de l'actualització?
  \item \textbf{(Interpretació)} L'IPC d'un any és del $0\%$ (els preus no han pujat de
        mitjana). Què vol dir per a una família? I si l'IPC fos \emph{negatiu}?
  \item \textbf{(Raonament)} Per què és important, per a un tècnic de serveis a la
        comunitat, saber interpretar l'IPC i les taxes de variació? Posa un exemple de
        com aquesta informació pot ajudar a defensar els drets d'una família. (La
        sessió següent veurem què passa quan el lloguer puja però l'ajut no.)
\end{exlist}
\end{exercicis}

% ---------------------------------------------------------------------
\fullsolucions{Sessió 10}
\begin{solucions}
\begin{sollist}
  \item Nova quota: $\num{700}\per 1,04=\num{728}\eur$. Augment mensual $\num{28}\eur$;
        augment anual $12\per\num{28}=\num{336}\eur$.
  \item Pujada del $6,5\%$ (índex $106,5$ sobre base $100$). El producte de
        $\num{20}\eur$ costa ara $\num{20}\per 1,065=\num{21,3}\eur$.
  \item $\dfrac{108,4}{100}=1,084$, és a dir un $+8,4\%$. Coincideix amb la taula.
  \item $\num{1100}\per 1,03=\num{1133}\eur$.
  \item Un IPC del $0\%$ vol dir que, de mitjana, els preus \textbf{no han pujat}
        respecte de l'any anterior: el poder de compra es manté. Un IPC
        \textbf{negatiu} (deflació) voldria dir que els preus han \textbf{baixat} de
        mitjana, de manera que amb els mateixos diners es podria comprar una mica més.
  \item Resposta oberta. Idea: entendre l'IPC permet comprovar si una actualització de
        lloguer és correcta, si una pensió o un ajut manté el poder de compra, o si una
        família està perdent capacitat de pagament. Exemple: detectar que el lloguer
        d'una família s'ha apujat \emph{per sobre} de l'IPC permetria reclamar-ho i
        defensar-ne els drets.
\end{sollist}
\end{solucions}
